El concepto de dataset refiere al conjunto de datos intencionalmente elegidos para ser utilizados en el entrenamiento, validación o prueba de un modelo de aprendizaje automático. Un dataset bien diseñado es crucial para el éxito de cualquier proyecto de aprendizaje automático, ya que la calidad y representatividad de los datos pueden influir significativamente en el rendimiento del modelo.
A modo de ejemplo, en el contexto de este proyecto se han utilizado dos datasets diferentes para los procesos de entrenamiento-validación y evaluación. La intención detrás de esta elección es garantizar que el modelo se entrene y valide con datos que reflejen una amplia variedad de casos, mientras que la evaluación se realiza con un conjunto de datos completamente independiente para medir la capacidad de generalización del modelo. Esta estrategia ayuda a evitar problemas como el overfitting, donde el modelo se ajusta demasiado a los datos de entrenamiento y no generaliza bien a nuevos datos.